%% The cost model behind the filter's benefit: what one write costs the cluster
%% under the original protocol, that no term of it depends on cache state, and
%% how the filter makes it. The bracket check uses results/oracle_ceiling.csv
%% (sender term = baseline_snoop - oracle_snoop, over four DCUs) and
%% results/invalidation_use.csv (receiver bound = broadcasts, over four DCUs);
%% the saving is baseline - oracle cycles of the slowest PE.
\subsection{Why suppression pays}
\label{sec:model}

Safety says the broadcast may be dropped; this says what dropping it is worth,
from the protocol rather than from the measurement.

Let core $i$ issue a write $w$ that reaches stage~1 of its DCU at cycle $a(w)$.
The original protocol grants it at the first cycle in which the arbiter selects
it \emph{and} every other DCU is ready to accept an invalidation,
\begin{equation}
g(w) = \min\Big\{t \ge a(w) : \mathrm{arb}(t) = i \;\wedge\; \textstyle\bigwedge_{j \ne i} r_j(t)\Big\}.
\label{eq:grant}
\end{equation}
A DCU is ready, $r_j(t) = 1$, only while its stage~1 is not already holding an
invalidation, and an invalidation leaves stage~1 only when stage~2 is free --
which it is not for the full memory latency of any miss or write-through in
flight. On grant, one invalidation is injected into stage~1 of each of the other
$P{-}1$ DCUs, where it takes precedence over the local core: a request already
there is displaced to a holding slot, one the core wanted to issue is refused,
and either way the local pipeline is shifted by one cycle. The cost of $w$ to
the cluster is therefore
\begin{equation}
C(w) = \big(g(w) - a(w)\big) + \sum_{j \ne i} \rho_j\big(g(w)\big),
\label{eq:cost}
\end{equation}
with $\rho_j(t) \in \{0, 1\}$ one when DCU $j$ had a request to shift and zero
when it was idle.

Two things follow. The conjunction in (\ref{eq:grant}) couples every pipeline
in the cluster: a store on core $i$ waits for a load miss on core $j$ to
complete, whether or not the two touch a common line. That is why the stall
share of Section~\ref{sec:measure} grows with memory latency -- the conjunction
is held false for the length of each miss. And neither term of (\ref{eq:cost})
depends on what any cache holds; the protocol charges $C(w)$ for every write
alike.

The filter makes it depend. With the sharer set $S(w)$ read from the mirror,
the grant condition becomes $\mathrm{arb}(t) = i \wedge \big(S(w) = \emptyset
\vee \bigwedge\nolimits_{j \ne i} r_j(t)\big)$ and invalidations are injected only when
$S(w) \ne \emptyset$. For a write with no sharers both terms of (\ref{eq:cost})
vanish, leaving the arbitration wait alone, which is non-zero only in a cycle
where two cores present writes at once. Summed over a run in which every write
has an empty sharer set -- which Section~\ref{sec:measure} measured to be the
case on all four kernels -- the saving is the whole coherence overhead less
arbitration collisions: exactly the oracle, which is why the filter reaches it
cycle for cycle rather than approximately.

Both terms are observable, so the model can be checked rather than fitted. The
sender term is the snoop-stall counter; the receiver term is at most one cycle
per invalidation delivered; so the two bracket the saving. Per DCU, at a memory
latency of two, the measured saving sits inside that bracket on every kernel:
matmul $120{,}539 \le 243{,}100 \le 326{,}373$ cycles, conv2d $5{,}185 \le
13{,}795 \le 39{,}980$, FFT $892 \le 4{,}865 \le 16{,}312$, memcpy $1{,}792 \le
2{,}302 \le 7{,}946$. The share of the receiver bound that materialises -- 0.60,
0.25, 0.26 and 0.08 -- is highest on matmul, whose pipeline is busy with in-cache
reads when the invalidations land, and lowest on memcpy, where the displaced
request was waiting on a write-through anyway.
